\documentclass{article}

\usepackage[final]{neurips_2026}

\usepackage[utf8]{inputenc}
\usepackage[T1]{fontenc}
\usepackage{hyperref}
\usepackage{url}
\usepackage{booktabs}
\usepackage{amsfonts}
\usepackage{nicefrac}
\usepackage{microtype}
\usepackage{xcolor}
\usepackage{graphicx}
\usepackage{subcaption}
\usepackage{amsmath}

\title{Early Warning System for Telco Customer Churn:\\
Predicting At-Risk Customers and Prioritizing Retention}

\author{%
  Pranjal Patel \\
  \texttt{pbp88} \\
  \And
  Shlok Bohra \\
  \texttt{sb2396} \\
  \And
  Romaan Roshan \\
  \texttt{rr1346} \\
}


\begin{document}

\maketitle

\begin{abstract}
Customer churn poses a significant financial challenge for telecommunications providers,
where retaining an existing customer is considerably less costly than acquiring a new one.
This paper presents an early warning system for predicting customer churn using the IBM
Telco Customer Churn dataset comprising 7,043 customers and 33 features.
We develop a rigorous end-to-end machine learning pipeline that addresses target leakage
removal, class imbalance, multicollinearity, and hyperparameter optimization.
Three classifiers are compared---logistic regression, random forest, and gradient
boosting---each trained within an \texttt{imblearn} pipeline using stratified
five-fold cross-validation and tuned to maximize recall.
Feature selection combines chi-squared tests, point-biserial correlation, and variance
inflation factor analysis to retain 27 statistically significant predictors.
The random forest achieves the highest recall (0.826) on the held-out test set, while
gradient boosting achieves the best AUC-ROC (0.853).
SHAP analysis identifies contract type, customer tenure, and internet service type
as the dominant churn drivers.
Cost-based threshold tuning further aligns model outputs with the asymmetric
business cost of missing a churner versus issuing a false retention offer.
Results demonstrate that interpretable, deployment-ready churn prediction is achievable
with standard machine learning methods when the pipeline is carefully constructed.
\end{abstract}


\section{Introduction}

Customer churn---the rate at which subscribers cancel or fail to renew their service---is
one of the most consequential operational challenges facing subscription-based businesses.
In the telecommunications sector, where service commoditization and low switching costs
intensify competitive pressure, retaining existing customers is estimated to be five to
seven times less expensive than acquiring new ones~\citep{vandenPoel2004}.
Even a modest reduction in monthly churn can translate to millions of dollars in preserved
revenue, making proactive churn prediction a high-value application of data science.

The core research question of this work is: \emph{can we accurately predict which
telecommunications customers are likely to churn, and which features most strongly drive
that decision?} We frame churn prediction as a binary classification problem and develop
an end-to-end pipeline on the publicly available IBM Telco Customer Churn
dataset~\citep{ibmTelco2019}.

Churn prediction is inherently asymmetric in cost. A \emph{false negative}---failing to
identify a customer who subsequently leaves---results in lost customer lifetime value.
A \emph{false positive}---flagging a loyal customer for a retention intervention---results
in a wasted discount or support call, which is comparatively inexpensive. This asymmetry
motivates prioritizing \emph{recall} over accuracy throughout our pipeline, and drives
our use of cost-based threshold tuning to translate predicted probabilities into
deployment-ready binary decisions.

Our contributions are as follows:
\begin{itemize}
  \item A complete preprocessing pipeline that explicitly handles target leakage, missing
        values, multicollinearity, and class imbalance.
  \item A rigorous feature selection stage combining chi-squared tests, point-biserial
        correlation, and variance inflation factor (VIF) analysis, each with documented
        decision thresholds.
  \item Comparison of logistic regression, random forest, and gradient boosting under
        stratified cross-validation with hyperparameter tuning.
  \item SHAP-based model interpretation connecting feature impacts to actionable business
        recommendations.
  \item Cost-based threshold optimization grounded in the business cost asymmetry between
        false negatives and false positives.
\end{itemize}


\section{Related Work}

\paragraph{Churn prediction in telecommunications.}
Early work on churn prediction relied heavily on survival analysis and logistic regression
to model subscriber attrition~\citep{vandenPoel2004}. \citet{verbeke2012} conducted a
comprehensive benchmark of classification methods for telecom churn, finding that
ensemble methods and profit-driven evaluation metrics outperform classical approaches when
the cost asymmetry between churn types is explicitly modeled. \citet{coussement2008}
demonstrated the utility of support vector machines for subscription churn but noted the
importance of feature engineering specific to the billing and contract domain.

\paragraph{Ensemble methods.}
Random forests~\citep{breiman2001} and gradient boosting machines~\citep{friedman2001}
have become the dominant approaches in tabular classification benchmarks. Random forests
reduce variance through bootstrap aggregation of decorrelated trees, while gradient
boosting reduces bias by sequentially fitting residuals. Both methods provide
feature importance scores as a natural by-product of training, which we leverage for
interpretation alongside SHAP values.

\paragraph{Class imbalance.}
Churn datasets are structurally imbalanced: churners typically represent 20--30\% of
customers. Standard accuracy maximization silently favors predicting the majority class.
\citet{chawla2002} introduced SMOTE (Synthetic Minority Oversampling Technique), which
generates synthetic minority samples by interpolating between existing ones.
\citet{heGarcia2009} provide a comprehensive survey of imbalance learning strategies,
including cost-sensitive learning via class weighting, which we compare empirically
against SMOTE.

\paragraph{Model interpretability.}
\citet{lundberg2017} introduced SHAP (SHapley Additive exPlanations), a unified
framework rooted in cooperative game theory that assigns each feature a contribution
value consistent with local model behavior. Unlike permutation importance or Gini
impurity, SHAP values are both locally faithful and globally consistent, making them
the preferred tool for explaining individual predictions in deployment
settings~\citep{molnar2020}.

\paragraph{Evaluation under imbalance.}
\citet{fawcett2006} established ROC analysis as the standard for comparing classifiers
across operating conditions. Under class imbalance, precision-recall curves are
more informative than ROC curves because they focus on classifier behavior over the
minority class, avoiding the optimistic distortion introduced by the large number of
true negatives~\citep{davis2006}.


\section{Method}

\subsection{Dataset}

We use the IBM Telco Customer Churn dataset~\citep{ibmTelco2019}, which contains
records for 7,043 customers of a fictional California telecommunications provider.
Each record includes 33 features spanning four categories: demographic attributes
(gender, senior citizen status, partner, dependents), account characteristics
(tenure, contract type, paperless billing, payment method), subscribed services
(phone, internet, online security, streaming), and billing (monthly charges, total charges).
The binary target variable \texttt{Churn Value} indicates whether the customer left
within the last month. The overall churn rate is 26.5\% (1,869 churners),
yielding a class imbalance ratio of approximately 2.8:1.

\subsection{Preprocessing}

\paragraph{Leakage removal.}
Four columns were removed prior to any modeling due to direct or proxy leakage of the
target: \texttt{Churn Score} (a precomputed churn propensity score), \texttt{CLTV}
(customer lifetime value computed post-churn), \texttt{Churn Reason} (free-text
exit survey field), and \texttt{Churn Label} (a string duplicate of the target).
Including any of these would constitute data leakage and invalidate all downstream results.

\paragraph{Non-informative columns.}
Nine geographic and identifier columns were removed: \texttt{CustomerID},
\texttt{Count}, \texttt{Country}, \texttt{State}, \texttt{City}, \texttt{Zip Code},
\texttt{Lat Long}, \texttt{Latitude}, and \texttt{Longitude}. These carry no
predictive signal for individual churn behavior.

\paragraph{Type correction.}
\texttt{Total Charges} is stored as a string column containing approximately 11
blank-space entries corresponding to customers with zero tenure (i.e., customers
enrolled mid-period who have not yet been billed). These entries were coerced to
numeric and imputed with zero, consistent with the interpretation that no charges
have accrued.

\paragraph{Encoding.}
All 16 remaining categorical columns were one-hot encoded with \texttt{drop\_first=True}
to avoid the dummy variable trap. The resulting design matrix has 30 features prior
to selection.

\subsection{Feature Selection}

Three complementary statistical tests were applied, each with a documented decision rule:

\textbf{Chi-squared test} (binary features). For each of the 27 binary columns produced
by one-hot encoding, a $\chi^2$ independence test was conducted against the churn label
using training data only. Features with $p \geq 0.05$ were dropped. Three features were
removed: \texttt{Gender\_Male}, \texttt{Phone Service\_Yes}, and
\texttt{Multiple Lines\_No phone service}.

\textbf{Point-biserial correlation} (continuous features). For each of the three
continuous features---\texttt{Tenure Months}, \texttt{Monthly Charges},
\texttt{Total Charges}---the point-biserial correlation with the binary churn target
was computed. All three exceeded the significance threshold ($p < 0.05$) and were retained.

\textbf{Variance inflation factor} (multicollinearity). VIF was computed for all retained
continuous features. The maximum observed VIF was 8.21 for \texttt{Total Charges},
below the standard threshold of 10. No features were dropped at this stage.

The final feature set consists of 27 predictors.

\subsection{Class Imbalance Handling}

Two strategies were evaluated via five-fold stratified cross-validation on logistic
regression as a representative model:

\begin{enumerate}
  \item \textbf{Class weighting}: \texttt{class\_weight=`balanced'} scales the loss
        contribution of minority class samples inversely proportional to class frequency.
  \item \textbf{SMOTE}: Synthetic minority samples are generated by interpolating
        between $k$-nearest neighbors in the minority class. Critically, SMOTE is
        applied \emph{inside} the cross-validation pipeline using \texttt{imblearn.pipeline},
        ensuring that synthetic samples from a given validation fold's neighborhood
        never contaminate that fold's training data.
\end{enumerate}

Both strategies produced nearly identical cross-validated recall (0.814 vs.\ 0.801)
and F1 scores (0.642 vs.\ 0.645). We selected class weighting for all final models
because it avoids synthetic data generation and introduces no additional hyperparameters.

\subsection{Model Training}

All models were wrapped in an \texttt{imblearn.Pipeline} to enforce that preprocessing
(scaling for logistic regression) is fitted on training folds only. Hyperparameter
search used \texttt{GridSearchCV} with \texttt{StratifiedKFold}($k=5$), optimizing recall.

\textbf{Logistic regression.} Regularization strength $C \in \{0.01, 0.1, 1, 10\}$
was searched. Features were standardized via \texttt{StandardScaler} inside the pipeline.
Best configuration: $C = 1$.

\textbf{Random forest.} Number of estimators $\in \{100, 200\}$ and maximum depth
$\in \{4, 8, \text{None}\}$ were searched. Best configuration: 100 estimators,
maximum depth 4.

\textbf{Gradient boosting.} Learning rate $\in \{0.05, 0.1\}$, number of estimators
$\in \{100, 200\}$, and maximum depth $\in \{3, 5\}$ were searched. Best configuration:
200 estimators, learning rate 0.05, maximum depth 3.

An 80/20 stratified train-test split was used throughout. The test set was held out
entirely until final evaluation; no decisions were made based on test set performance.

\subsection{Threshold Optimization}

By default, models classify a customer as a churner if their predicted probability
exceeds 0.5. Given the asymmetric cost structure, we optimize the threshold $\tau^*$ by
minimizing:
\begin{equation}
  \mathcal{C}(\tau) = c_{\text{FP}} \cdot (1 - \text{Precision}(\tau))
                    + c_{\text{FN}} \cdot (1 - \text{Recall}(\tau))
\end{equation}
where $c_{\text{FP}}$ and $c_{\text{FN}}$ are the relative costs of false positives and
false negatives respectively. We use a ratio of $c_{\text{FN}} : c_{\text{FP}} = 10:1$,
reflecting the practical asymmetry between lost customer lifetime value and a wasted
retention discount. The optimal threshold is identified by sweeping over the precision-recall
curve of the best-performing model.


\section{Evaluation}

\subsection{Dummy Baseline}

Prior to evaluating learned models, a \texttt{DummyClassifier} predicting the majority
class for every input was established as a performance floor. This baseline achieves
73.4\% accuracy but 0\% recall on the churn class, confirming that accuracy is a
misleading metric under class imbalance and that any useful model must substantially
outperform this floor on recall.

\subsection{Model Performance}

Table~\ref{tab:results} reports precision, recall, F1 score, and AUC-ROC for all
three models on the held-out test set at the default threshold of 0.5.

\begin{table}[h]
  \caption{Test-set performance of all three models at threshold 0.5.
           Random forest achieves the highest recall; gradient boosting achieves
           the highest AUC-ROC.}
  \label{tab:results}
  \centering
  \begin{tabular}{lcccc}
    \toprule
    Model & Precision & Recall & F1 & AUC-ROC \\
    \midrule
    Dummy Classifier        & 0.000 & 0.000 & 0.000 & 0.500 \\
    Logistic Regression     & 0.511 & 0.783 & 0.619 & 0.849 \\
    Random Forest           & 0.512 & \textbf{0.826} & \textbf{0.632} & 0.846 \\
    Gradient Boosting       & \textbf{0.656} & 0.540 & 0.592 & \textbf{0.853} \\
    \bottomrule
  \end{tabular}
\end{table}

All three learned models achieve substantially higher recall than the dummy baseline.
The random forest is the best model for deployment scenarios prioritizing recall,
catching 82.6\% of actual churners on the test set. Gradient boosting, while lower in
recall at the default threshold, achieves the best overall discriminative ability
(AUC-ROC 0.853) and the highest precision (0.656), making it preferable when
intervention budgets are constrained.

\begin{figure}[h]
  \centering
  \begin{subfigure}[b]{0.48\linewidth}
    \includegraphics[width=\linewidth]{roc_curves.png}
    \caption{ROC curves for all three models. Gradient boosting marginally leads
             in AUC-ROC (0.853).}
    \label{fig:roc}
  \end{subfigure}
  \hfill
  \begin{subfigure}[b]{0.48\linewidth}
    \includegraphics[width=\linewidth]{pr_curves.png}
    \caption{Precision-Recall curves. All models substantially exceed the random
             baseline (dashed line at 0.27).}
    \label{fig:pr}
  \end{subfigure}
  \caption{Discrimination performance on the held-out test set.}
  \label{fig:curves}
\end{figure}

\subsection{Overfitting Analysis}

To verify generalization, we compare training-set recall against mean cross-validated
recall for each model. The gaps are 0.002 (logistic regression), 0.014 (random forest),
and 0.053 (gradient boosting), all below the 0.08 threshold we set as a warning criterion.
No model exhibits meaningful overfitting, consistent with the regularization introduced
by depth constraints and class weighting.

\subsection{Confusion Matrices}

Figure~\ref{fig:cm} shows the confusion matrices for all three models at the default
threshold. The random forest misses only 65 churners (false negatives) compared to 81
for logistic regression, while gradient boosting is considerably more conservative,
missing 172 churners but flagging fewer loyal customers incorrectly.

\begin{figure}[h]
  \centering
  \includegraphics[width=0.95\linewidth]{confusion_matrices.png}
  \caption{Confusion matrices on the test set (threshold = 0.5). Random forest
           catches the most churners (309 true positives, 65 false negatives).}
  \label{fig:cm}
\end{figure}

\subsection{Threshold Optimization}

Applying cost-based threshold optimization to the gradient boosting model (highest
AUC-ROC) with cost ratio $c_{\text{FN}}:c_{\text{FP}} = 10:1$ yields an optimal
threshold of 0.023. At this threshold, recall increases to 1.000 and precision
decreases to 0.311. While a threshold this low is aggressive, it illustrates the
direct connection between the business cost ratio and the deployment operating point:
organizations with high customer lifetime value and low retention offer costs should
operate at correspondingly lower thresholds. Figure~\ref{fig:threshold} shows the
cost curve and the precision-recall trade-off across all thresholds.

\begin{figure}[h]
  \centering
  \includegraphics[width=0.75\linewidth]{threshold_tuning.png}
  \caption{Cost-based threshold optimization for gradient boosting ($c_{\text{FN}}=10$,
           $c_{\text{FP}}=1$). The red dashed line marks the optimal threshold (0.023)
           that minimizes total expected cost.}
  \label{fig:threshold}
\end{figure}


\section{Discussion}

\subsection{Feature Importance and SHAP Analysis}

SHAP values were computed for the random forest using \texttt{TreeExplainer}, which
exactly decomposes each prediction into per-feature contributions consistent with
Shapley values from cooperative game theory. Figure~\ref{fig:shap} shows the mean
absolute SHAP impact across the test set.

\begin{figure}[h]
  \centering
  \includegraphics[width=0.75\linewidth]{shap_bar.png}
  \caption{Mean absolute SHAP values for the random forest. Contract type and
           customer tenure are the two most impactful predictors of churn.}
  \label{fig:shap}
\end{figure}

The five most impactful features are:

\textbf{Contract type (Two-year contract).} Being on a two-year contract is the
strongest predictor of \emph{reduced} churn risk. Month-to-month customers, lacking
contractual lock-in, exhibit the highest attrition rates. This is consistent with
findings in~\citet{verbeke2012}.

\textbf{Tenure months.} Longer tenure strongly correlates with loyalty. The
point-biserial correlation of $-0.346$ confirms that the relationship is inverse:
newer customers are significantly more likely to churn. This suggests a critical
intervention window in the first 12 months.

\textbf{Internet service (Fiber Optic).} Fiber Optic subscribers churn at
disproportionately high rates despite paying premium prices, suggesting unmet
quality expectations rather than price sensitivity.

\textbf{Dependents.} Customers with dependents are less likely to churn, potentially
because service switching involves coordinating across a household.

\textbf{Electronic check payment.} This payment method is associated with elevated
churn risk, possibly as a proxy for lower service engagement or financial instability.

\subsection{Business Recommendations}

Based on SHAP findings, we derive five specific, quantifiable recommendations:

\begin{enumerate}
  \item \textbf{Early tenure intervention.} Offer discounted annual contracts to
        month-to-month customers at month 6, before the most common churn window.
        Customers on month-to-month contracts with tenure under 12 months represent
        the modal churner profile in this dataset.

  \item \textbf{Add-on bundling for high-charge customers.} Customers paying above
        \$65/month without \texttt{Online Security} or \texttt{Tech Support} show
        elevated churn risk. Bundling these services into a loyalty tier can simultaneously
        reduce churn and increase average revenue per user.

  \item \textbf{Deployment threshold calibration.} Replace the default threshold of 0.5
        with the cost-optimized threshold derived from the organization's estimated
        $c_{\text{FN}}:c_{\text{FP}}$ ratio. At a 10:1 ratio, the threshold drops
        substantially, prioritizing recall over precision.

  \item \textbf{Fiber Optic quality audit.} Before deploying price-based retention
        incentives to Fiber Optic subscribers, conduct a service quality survey.
        High churn among premium subscribers is a service quality signal, not
        a pricing signal.

  \item \textbf{Senior citizen support program.} Senior citizens exhibit
        above-average churn rates, likely reflecting difficulty with self-service
        channels or billing complexity. A dedicated support tier for this segment
        can reduce churn without requiring discounts.
\end{enumerate}

\subsection{Model Selection Trade-offs}

Each model occupies a distinct position in the interpretability-performance trade-off space.
Logistic regression provides the most transparent predictions via signed coefficients with
confidence intervals, making it preferable for regulated environments where prediction
explanations must be provided to customers. Random forest maximizes recall at the default
threshold, making it the best default deployment choice for high-volume outreach programs.
Gradient boosting maximizes AUC-ROC and precision, making it optimal for constrained
retention budgets where false positives are costly.


\section{Limitations}

\textbf{Dataset currency and scope.} The IBM Telco dataset is a synthetic educational
sample rather than a live customer database. It does not reflect current market
dynamics, carrier-specific customer behavior, or temporal effects such as pricing
changes or competitive events. Models trained on this dataset should not be deployed
without validation on proprietary operational data.

\textbf{Static modeling.} Our models capture a single snapshot of customer behavior.
Churn drivers evolve over time as pricing, technology, and competition change.
Production deployment would require scheduled retraining and monitoring for
covariate shift.

\textbf{No causal identification.} High SHAP impact does not imply causation.
The association between contract type and churn, for example, may partially reflect
selection effects (i.e., customers who intend to leave are less likely to sign
long-term contracts in the first place). Retention strategies based on our findings
should be validated with randomized A/B experiments before large-scale rollout.

\textbf{Cost ratio sensitivity.} The optimal threshold depends directly on the
assumed $c_{\text{FN}}:c_{\text{FP}}$ ratio. We used an illustrative ratio of 10:1;
in practice, this ratio should be estimated from customer lifetime value models
and the cost structure of the retention intervention program.

\textbf{Geographic features excluded.} City and zip code features were dropped as
non-informative for a synthetic dataset. In real deployment, geographic features
may capture significant variation in churn behavior due to regional competition
or infrastructure quality.


\begin{ack}
This project was completed as part of CS 439. No external funding was received.
The IBM Telco Customer Churn dataset is publicly available on Kaggle under a
community data license.
\end{ack}


\section*{References}

{
\small

[1] Breiman, L. (2001). Random forests. \textit{Machine Learning}, 45(1), 5--32.

[2] Chawla, N.V., Bowyer, K.W., Hall, L.O., \& Kegelmeyer, W.P. (2002). SMOTE:
Synthetic minority over-sampling technique. \textit{Journal of Artificial Intelligence
Research}, 16, 321--357.

[3] Coussement, K. \& Van den Poel, D. (2008). Churn prediction in subscription
services using support vector machines. \textit{Expert Systems with Applications},
34(1), 313--327.

[4] Davis, J. \& Goadrich, M. (2006). The relationship between Precision-Recall and
ROC curves. In \textit{Proceedings of the 23rd ICML}, pp.\ 233--240.

[5] Fawcett, T. (2006). An introduction to ROC analysis. \textit{Pattern Recognition
Letters}, 27(8), 861--874.

[6] Friedman, J.H. (2001). Greedy function approximation: a gradient boosting machine.
\textit{Annals of Statistics}, 29(5), 1189--1232.

[7] He, H. \& Garcia, E.A. (2009). Learning from imbalanced data. \textit{IEEE
Transactions on Knowledge and Data Engineering}, 21(9), 1263--1284.

[8] IBM Corporation. (2019). Telco Customer Churn Dataset. Kaggle.
\url{https://www.kaggle.com/datasets/blastchar/telco-customer-churn}

[9] Lundberg, S.M. \& Lee, S.I. (2017). A unified approach to interpreting model
predictions. In \textit{Advances in Neural Information Processing Systems}, 30.

[10] Molnar, C. (2020). \textit{Interpretable Machine Learning}. Lulu.com.

[11] Van den Poel, D. \& Larivi\`{e}re, B. (2004). Customer attrition analysis for
financial services using proportional hazard models. \textit{European Journal of
Operational Research}, 157(1), 196--217.

[12] Verbeke, W., Dejaeger, K., Martens, D., Hur, J., \& Baesens, B. (2012). New
insights into churn prediction in the telecommunication sector: A profit driven data
mining approach. \textit{European Journal of Operational Research}, 218(1), 211--229.

}


\appendix

\section{Additional Results}

Table~\ref{tab:featsel} summarizes the feature selection decisions made at each stage
of the selection pipeline.

\begin{table}[h]
  \caption{Feature selection summary. 27 of 30 encoded features were retained.}
  \label{tab:featsel}
  \centering
  \begin{tabular}{lccc}
    \toprule
    Stage & Test & Features evaluated & Features dropped \\
    \midrule
    Chi-squared      & $\chi^2$, $p < 0.05$     & 27 binary   & 3 \\
    Point-biserial   & $r_{pb}$, $p < 0.05$     & 3 continuous & 0 \\
    VIF              & VIF $\leq 10$            & 3 continuous & 0 \\
    \midrule
    \textbf{Final}   &                           & \textbf{27} & \textbf{3} \\
    \bottomrule
  \end{tabular}
\end{table}

The three features removed by chi-squared testing were \texttt{Gender\_Male}
($p = 0.481$), \texttt{Phone Service\_Yes} ($p = 0.481$), and
\texttt{Multiple Lines\_No phone service} ($p = 0.481$). The maximum VIF observed
among continuous features was 8.21 (\texttt{Total Charges}), below the
multicollinearity threshold of 10.

\end{document}
